\section*{§1. Sets}

\textbf{Definition 1.} A \textit{set} is a collection of objects. These objects are called the \textit{elements} of the set.

We will usually denote sets by capital letters $A, B, C, \dots$ and elements by small letters $a, b, c, \dots$. If $a$ is an element of the set $A$, we write $a \in A$ (read "$a$ belongs to $A$"). If $a$ is not an element of $A$, we write $a \notin A$ (read "$a$ does not belong to $A$").

A set can be defined by listing all its elements between braces. For example, the set $A$ consisting of the numbers 1, 2, and 3 is written as
\[
A = \{1, 2, 3\}.
\]
The order in which the elements are listed is irrelevant. Thus, the sets $\{1, 2, 3\}$ and $\{3, 1, 2\}$ are the same. Also, if an element is listed more than once, it is counted only once. Thus, the sets $\{1, 2, 3\}$ and $\{1, 2, 3, 2\}$ are the same.

A set can also be defined by specifying a property that all its elements satisfy. For example, the set of all real numbers $x$ such that $x^2 < 2$ is written as
\[
\{x \in \mathbb{R} \mid x^2 < 2\}.
\]
In general, if $P(x)$ is a property that $x$ may or may not satisfy, then the set of all $x$ in $A$ such that $P(x)$ is true is written as
\[
\{x \in A \mid P(x)\}.
\]
The symbol "|" is read "such that".

\textbf{Definition 2.} The set that contains no elements is called the \textit{empty set} and is denoted by $\emptyset$.

\textbf{Definition 3.} If every element of a set $A$ is also an element of a set $B$, then $A$ is called a \textit{subset} of $B$, and we write $A \subset B$ (read "$A$ is contained in $B$"). If $A \subset B$ and $A \neq B$, then $A$ is called a \textit{proper subset} of $B$.

It follows from the definition that the empty set is a subset of every set. Also, every set is a subset of itself.

\textbf{Definition 4.} If $A \subset B$ and $B \subset A$, then $A$ and $B$ are said to be \textit{equal}, and we write $A = B$.

Thus, two sets are equal if and only if they have the same elements.